\documentclass[11pt]{article}
\usepackage[margin=1in]{geometry}
\usepackage{times}
\usepackage{hyperref}
\hypersetup{colorlinks=true, linkcolor=black, urlcolor=blue, citecolor=black}
\title{Same Inventor, Same Substitution: Bruce DePalma's N-Machine Homopolar Generator}
\author{Flyxion \\ \small Independent Researcher}
\date{}
\begin{document}
\maketitle

\begin{abstract}
Bruce DePalma, in addition to the spinning-ball gravity experiments addressed elsewhere in this series, promoted from the late 1970s a homopolar generator design he called the N-machine, claiming that a specific configuration of a magnetized rotating conductor and unusual brush arrangement produced electrical output exceeding the mechanical energy required to spin it. This paper argues that DePalma's N-machine efficiency claims trace to the same category of mechanical-input-versus-electrical-input measurement substitution addressed in the Tewari Space Power Generator entry elsewhere in this series, that homopolar generators have well understood and entirely conventional efficiency characteristics independent of DePalma's specific claimed innovation, and that no independent, peer-reviewed replication using calibrated mechanical torque instrumentation has confirmed the claimed over-unity output in the decades since the design was first promoted.
\end{abstract}

\section{Introduction}

Bruce DePalma, following his spinning-ball gravity experiments discussed separately in this series, promoted from the late 1970s a homopolar generator design, a rotating magnetized disc or cylinder with brush contacts drawing current radially, which he called the N-machine, claiming that specific magnet and brush configurations allowed the device to produce electrical output exceeding the mechanical power required to maintain its rotation, framing this as related to his broader theoretical claims about rotation representing an unrecognized form of energy.

\section{Homopolar Generators Are Conventional but Have Documented Quirks}

The homopolar generator is a genuinely old and well understood class of electrical machine, first demonstrated by Faraday, that produces direct current through the motion of a conductor across a magnetic field without the alternating polarity reversal of conventional generators, and is known in conventional electrical engineering for a specific quirk relevant to this claim: a homopolar generator's magnetic drag torque, the resistance a rotor feels as current is drawn, can be comparatively small relative to a conventional generator's for certain geometries and load conditions, a real, entirely conventional property that has occasionally been mistaken, by researchers not fully accounting for it, as evidence the device is not obeying the usual electromechanical energy balance, when a complete accounting of internal resistive losses, brush contact losses, which are typically substantial in homopolar designs due to the sliding brush contacts, and the true total mechanical input restores the expected balance.

\section{The Same Measurement Substitution Addressed Elsewhere in This Series}

As with the Tewari Space Power Generator, published descriptions and demonstrations of DePalma's N-machine efficiency claims have generally relied on comparing the electrical input to the motor driving the generator, rather than a direct, independently instrumented measurement of the mechanical torque and rotational speed actually delivered to the generator's shaft, a substitution that systematically understates true mechanical input power by the driving motor's own conversion inefficiency and can produce an apparent generator efficiency exceeding one hundred percent even when the generator itself, examined on a properly instrumented torque and speed basis, performs within or below conventional expectations once brush contact losses, often substantial for homopolar designs specifically, are properly accounted for.

\section{Why DePalma's Two Separate Claims in This Series Share a Common Pattern}

The spinning-ball experiment addressed elsewhere in this series and the N-machine addressed here are physically distinct claims, one concerning gravitational mass and one concerning generator efficiency, but both share a common evidentiary weakness: a comparison between two quantities, launch velocity in one case and mechanical shaft power in the other, that was not independently and directly measured but was instead inferred from a related, more easily measured proxy, spring mechanism nominal specification in one case and motor electrical input in the other, that does not reliably track the actual quantity the underlying physical claim depends on, a pattern worth noting as a recurring methodological signature across DePalma's separate claimed discoveries rather than as two unrelated coincidences.

\section{What Rigorous Replication Has Found}

Independent evaluations of homopolar generator designs, including N-machine-inspired configurations, using calibrated torque sensors and direct mechanical power measurement rather than motor electrical input as a proxy, have consistently found electrical output within the bounds conventional electromagnetic theory predicts once brush contact resistance and other conventional losses are included, with no peer-reviewed publication in the decades since DePalma's original claims confirming a genuine over-unity result under this more rigorous measurement standard.

\section{The Entropic Gravity Framing}

The N-machine's claim concerns electrical generator efficiency rather than gravitational coupling directly, and is included in this series as a comparative case alongside the Tewari generator and other over-unity claims addressed elsewhere here, illustrating how the same specific measurement substitution, motor electrical input mistaken for mechanical shaft power, recurs across independently promoted devices spanning different inventors, decades, and claimed underlying mechanisms.

\section{Objections}

Supporters argue that DePalma's specific brush geometry and magnet configuration represent a genuine, underexplored variant of homopolar generator design deserving independent engineering study on its own terms rather than dismissal by association with the measurement substitution pattern. Independent engineering study of novel homopolar generator geometries is a legitimate and ongoing area of electrical machine design research, and this paper's critique is not that homopolar generators cannot be studied or improved, but specifically that the over-unity efficiency claims made for DePalma's particular design have not been confirmed by any replication using the calibrated, direct mechanical power measurement standard that would be needed to distinguish a genuine anomaly from the well documented measurement substitution addressed above.

\section{Conclusion}

The N-machine's claimed over-unity efficiency traces to the same motor-electrical-input-versus-mechanical-shaft-power measurement substitution documented in comparable over-unity generator claims elsewhere in this series, homopolar generators' genuine but conventional low-drag-torque properties do not, on proper accounting including brush contact losses, produce the claimed excess, and no independent, properly instrumented replication has confirmed the claim in the decades since it was first promoted.

\end{document}
